\subsection{Сравнение времени записи}

Эксперимент запускался последовательно для четырех моделей хранения:
\texttt{pg-jsonb}, \texttt{pg-normalized}, \texttt{mongo-nested} и
\texttt{mongo-normalized}. Для всех запусков использовались одинаковые
параметры: профиль данных \texttt{small}, \texttt{seed = 42}, ступени
параллелизма 1, 5, 10 и 25 клиентов, длительность каждой ступени 60 секунд.
Порядок запуска был фиксированным: сначала PostgreSQL JSONB, затем
нормализованная PostgreSQL, вложенная MongoDB и нормализованная MongoDB.

Запись проверялась двумя операциями: пакетной записью сложных событий
\texttt{WRITE\_COMPLEX\_EVENT\_BATCH} и потоковой записью телеметрии
\texttt{WRITE\_TELEMETRY\_STREAM}. В сценарии \texttt{write-heavy} на эти
операции приходится 95\% нагрузки, поэтому он лучше всего показывает цену
вставки данных в разные модели.

На рисунке~\ref{fig:dashboard-writeheavy} видно, что за прогон
\texttt{write-heavy} выполнено около 113 тыс. операций, обработано около
1.7 GiB данных, ошибки не зафиксированы. Самую высокую пропускную способность
показала вложенная MongoDB-модель: пик на графике пропускной способности по
моделям превышает 300 операций в секунду. PostgreSQL JSONB и нормализованная
PostgreSQL держатся ниже, но остаются устойчивыми. Нормализованная MongoDB
заметно отстает по пропускной способности и дает более длинные хвостовые
задержки.

\begin{figure}[H]
    \centering
    \includegraphics[
        width=\textwidth,
        height=1\textheight,
        keepaspectratio
    ]{dash-writeheavy.png}
    \caption{Панель Grafana для сценария write-heavy}
    \label{fig:dashboard-writeheavy}
\end{figure}

\subsection{Сравнение времени чтения}

Чтение проверялось тремя операциями: агрегацией активности объектов,
агрегацией технического состояния камер и чтением таймлайна инцидента.
Сценарий \texttt{analytics-heavy} отдает этим операциям 70\% нагрузки, поэтому
лучше показывает стоимость аналитических запросов.

За прогон \texttt{analytics-heavy} выполнено около 71 тыс. операций и
обработано около 812.6 MiB данных. На рисунке~\ref{fig:dashboard-analytics}
видно, что вложенная MongoDB-модель и нормализованная PostgreSQL дают
наибольшую пропускную способность в начале и середине прогона: примерно 140 и
90--100 операций в секунду соответственно. PostgreSQL JSONB показывает
хорошую скорость на отдельных ступенях, но задержки аналитических операций
растут при увеличении числа клиентов. Нормализованная MongoDB снова имеет
самую низкую пропускную способность, потому что аналитические запросы
вынуждены восстанавливать контекст через ссылки между коллекциями.

\begin{figure}[H]
    \centering
    \includegraphics[
        width=\textwidth,
        height=1\textheight,
        keepaspectratio
    ]{dash-analytics.png}
    \caption{Панель Grafana для сценария analytics-heavy}
    \label{fig:dashboard-analytics}
\end{figure}

\subsection{Сравнение моделей при разной нагрузке}

Сценарий \texttt{balanced} используется как смешанный режим: 70\% операций
связано с записью, 30\% -- с чтением и аналитикой. Рост нагрузки задавался
ступенями параллелизма. На каждой ступени loadgen запускал одинаковую смесь
операций, а Grafana показывала пропускную способность и процентили задержек
по модели, операции и числу клиентов.

\begin{figure}[H]
    \centering
    \includegraphics[
        width=\textwidth,
        height=1\textheight,
        keepaspectratio
    ]{dash-balanced.png}
    \caption{Панель Grafana после последовательного запуска моделей в сценарии balanced}
    \label{fig:dashboard-balanced}
\end{figure}

За один последовательный запуск было выполнено около 59 тыс. операций,
доля ошибок составила 0\%, объем обработанных данных -- около 994 MiB.
Пиковая пропускная способность различалась по моделям. На графике
пропускной способности по моделям вложенная MongoDB достигает примерно
150 операций в секунду, нормализованная PostgreSQL -- около 90 операций в
секунду, PostgreSQL JSONB -- около 60 операций в секунду. Нормализованная
MongoDB в этом запуске показала меньшую пропускную способность и более
длинные хвостовые задержки: p95 отдельных аналитических операций поднимается
примерно до 25--30 секунд.

\subsection{Анализ полученных результатов}

После всех экспериментов лучший результат показала вложенная MongoDB. В
сценарии \texttt{write-heavy} она дала наибольшую пропускную способность при
записи событий и телеметрии. Это связано с тем, что сложное событие
сохраняется как один документ вместе со всем своим контекстом,
а не раскладывается на несколько таблиц или коллекций.

В сценариях \texttt{analytics-heavy} и \texttt{balanced} вложенная MongoDB
также осталась лучшей по общему поведению. Для запросов системы
видеонаблюдения часто требуется получить событие вместе с территорией, зоной,
камерами и обнаруженными объектами. Во вложенной модели эти данные уже лежат
рядом, поэтому запросу нужно меньше соединений и переходов по ссылкам.

Нормализованная PostgreSQL показала себя лучше нормализованной MongoDB:
реляционные связи и типизированные поля дают более устойчивое выполнение
аналитических запросов. PostgreSQL JSONB заняла промежуточное положение:
вариативный \texttt{payload} удобно записывать, но агрегации по массивам
внутри JSONB дороже работы с обычными столбцами. Нормализованная MongoDB
оказалась самым слабым вариантом для аналитики: при чтении ей приходится
восстанавливать контекст через \texttt{\$lookup}.

Индексы были созданы во всех моделях по полям, которые участвуют в фильтрации
и связях: времени, территории, зоне, типу события, уровню значимости и
идентификаторам связанных сущностей. Поэтому результат нельзя объяснить тем,
что одна модель была проиндексирована, а другая нет. Разница возникла из-за
самой структуры хранения: вложенная MongoDB лучше соответствует тем запросам,
которые выполнялись в эксперименте.
